%%%%%%%%%%%%%%%%%%%%%%%%%%%%%%%%%%%%%%%%%%%%%%%%%%%%%%%%%%%%%%%%%%%%%%%%%%%%%
% Monster orbifold E_3-topological descent criterion.
%
% Separates the VOA input from the HT descent input: start with the
% class-G Leech HT model, assume anomaly-free finite-orbifold BV descent,
% and identify the descended boundary algebra with the FLM Moonshine
% module.
%%%%%%%%%%%%%%%%%%%%%%%%%%%%%%%%%%%%%%%%%%%%%%%%%%%%%%%%%%%%%%%%%%%%%%%%%%%%%

\chapter[Monster orbifold E\_3 descent]{Monster orbifold
$E_3$-topological descent via the Leech $\Z/2$ orbifold}
\label{ch:monster-chain-level-e3-top-platonic}

\paragraph{Orientation.}
The Moonshine module $V^\natural$ is obtained from the class-G Leech
lattice VOA $V_\Lambda$ by the FLM finite $\Z/2$ orbifold. That VOA
classification input is not an HT descent theorem. The HT statement
requires three further data: the abelian Leech holomorphic-topological
model, finite-orbifold BV descent for the lifted action, and a
trivialization of the Dijkgraaf-Witten obstruction
\[
 [\alpha_{\mathrm{orb}}]\in H^3(B\Z/2;\mathrm{U}(1))\cong \Z/2.
\]
In multiplicative normalization the trivial class is detected by
$\alpha_{\mathrm{orb}}(\sigma,\sigma,\sigma)=+1$. This sign is an
anomaly datum or a separately verified local computation, not a
consequence of the determinant slogan alone.

The chapter records the resulting descent criterion. It does not
upgrade the generic class-M DS-Hochschild statement beyond the
weight-completed ambient, and it does not assert that $V^\natural$ as
an abstract boundary VOA automatically carries a single-colour
original-complex $E_3$ structure.

%%%%%%%%%%%%%%%%%%%%%%%%%%%%%%%%%%%%%%%%%%%%%%%%%%%%%%%%%%%%%%%%%%%%%%%%%%%%%
\section{The theorem}
\label{sec:monster-chain-level-e3-top-statement}
%%%%%%%%%%%%%%%%%%%%%%%%%%%%%%%%%%%%%%%%%%%%%%%%%%%%%%%%%%%%%%%%%%%%%%%%%%%%%

\begin{theorem}[Monster orbifold $E_3$-topological descent criterion;
\ClaimStatusConditional]
\label{thm:monster-chain-level-e3-top}
Let $V_\Lambda$ be the Frenkel-Lepowsky-Meurman lattice vertex
operator algebra on the Leech lattice $\Lambda$, and
$\widetilde{\sigma}\colon V_\Lambda \to V_\Lambda$ the canonical FLM
lift of the $-1$ isometry (with cocycle normalisation
$c(\lambda) \equiv 1$). Let
\[
 V^\natural
 \;=\; V_\Lambda^+
 \,\oplus\, V_\Lambda^{\mathrm{tw},+}
\]
be the FLM Moonshine module, the $\widetilde{\sigma}$-invariant part
of the $\Z/2$-orbifold of $V_\Lambda$. Assume:
\begin{enumerate}[label=\textup{(\roman*)}]
\item the Leech lattice HT model $\cF_\Lambda$ carries the abelian
chain-level $E_3$-topological structure on the ordinary chain complex;
\item the lifted $\Z/2$ action admits finite-orbifold BV descent and
its obstruction class
$[\alpha_{\mathrm{orb}}]\in H^3(B\Z/2;\mathrm{U}(1))$ is trivial,
equivalently
$\alpha_{\mathrm{orb}}(\sigma,\sigma,\sigma)=+1$ after choosing a
normalized representative;
\item boundary restriction for the descended orbifold model identifies
the boundary chiral algebra with the FLM orbifold
$V_\Lambda^+\oplus V_\Lambda^{\mathrm{tw},+}$.
\end{enumerate}
Then the descended orbifold factorisation algebra
$\cF_\Lambda^{\mathrm{orb}}$ carries a chain-level $E_3$-topological
structure, and its boundary chiral algebra is $V^\natural$. This is a
finite-orbifold HT descent from the abelian Leech model, not an
automatic original-complex $E_3$ structure on the abstract Monster VOA.
\end{theorem}

\begin{proof}
FLM construct the boundary VOA:
$V_\Lambda^{\mathrm{tw}}$ has ground-state multiplicity $2^{12}$, and
FLM Chapters~8--10 prove the twisted Jacobi identity for
$V_\Lambda^+\oplus V_\Lambda^{\mathrm{tw},+}$. Dong--Mason quantum
Galois uniqueness fixes the twisted sector up to VOA isomorphism.
These are VOA statements.

Hypothesis~(i) supplies the $E_3$-topological operations before
gauging. Hypotheses~(ii)--(iii) say that the finite-orbifold BV
construction has no anomalous associator and that boundary restriction
identifies the descended boundary algebra with the FLM orbifold.
Costello--Gwilliam descent then transports the $E_3$ operations from
$\cF_\Lambda$ to $\cF_\Lambda^{\mathrm{orb}}$. The conclusion is
confined to the Leech finite-orbifold model and is compatible with the
generic class-M weight-completed scope of
Corollary~\ref*{cor:universal-holography-class-M}.
\end{proof}

\begin{remark}[Scope of the Monster exception]
\label{rem:monster-orbifold-scope}
The theorem gives a descended HT model beyond the generic class-M
statement only after the finite-orbifold BV hypotheses have been
supplied. The Monster is reached from class~G before the orbifold is
taken. The criterion requires:
\begin{itemize}
\item the anomaly trivialization for the Leech $(-1)$ orbifold;
\item the FLM twisted-sector chain algebra and Dong--Mason uniqueness
 (FLM Chapters~8--10 and Dong--Mason);
\item finite-orbifold BV descent of the abelian class-G
 $E_3$-topological structure.
\end{itemize}
External inputs remain external: the abelian-CS mechanism for
$V_\Lambda$ (Theorem~\ref*{thm:E3-topological-km}), the FLM twisted
Jacobi identity (FLM~1988), Dong--Mason uniqueness, and any local
verification of the BV anomaly datum. None is replaced by the formal
identity $\det(1-(-1)|_\Lambda)=2^{24}$.
\end{remark}

%%%%%%%%%%%%%%%%%%%%%%%%%%%%%%%%%%%%%%%%%%%%%%%%%%%%%%%%%%%%%%%%%%%%%%%%%%%%%
\section{The $\Z/2$ anomaly datum}
\label{sec:monster-anomaly-class-explicit}
%%%%%%%%%%%%%%%%%%%%%%%%%%%%%%%%%%%%%%%%%%%%%%%%%%%%%%%%%%%%%%%%%%%%%%%%%%%%%

The Dijkgraaf-Witten 3-cocycle group $H^3(B\Z/2;\mathrm{U}(1))$ is
generated by a single sign: either $\alpha(\sigma,\sigma,\sigma) = 1$
(trivial class $0$) or $\alpha(\sigma,\sigma,\sigma) = -1$
(non-trivial class $\omega$, corresponding to the fractional
Chern-Simons theory $\Z/2_2$ with quadratic Pontryagin action).
For the FLM lift the equality
\[
 \alpha_{\mathrm{orb}}(\sigma,\sigma,\sigma)=+1
\]
is the finite-orbifold BV datum needed by the descent theorem. The
numerical fact
\[
 \det(1-\sigma|_\Lambda)=\det(2\cdot\mathrm{id}_\Lambda)=2^{24}
\]
checks that the real determinant sign contributes $+1$; it does not
by itself compute the full group-cohomology obstruction.

\begin{proposition}[Leech $\Z/2$ anomaly criterion;
\ClaimStatusConditional]
\label{prop:monster-alpha-explicit-zero}
For the FLM $\Z/2$ orbifold of $V_\Lambda$, finite-orbifold BV
descent is unobstructed exactly when the chosen local anomaly datum
satisfies
\[
 \alpha_{\mathrm{orb}}(\sigma,\sigma,\sigma)=+1.
\]
Equivalently, $[\alpha_{\mathrm{orb}}]=0$ in
$H^3(B\Z/2;\mathrm{U}(1))\cong\Z/2$.
\end{proposition}

\begin{proof}
The group-cohomology statement is
$H^3(B\Z/2;\mathrm{U}(1))\cong\Z/2$. In normalized multiplicative
coordinates the two classes are represented by
$\alpha(\sigma,\sigma,\sigma)=+1$ and
$\alpha(\sigma,\sigma,\sigma)=-1$. The former is the trivial
Dijkgraaf-Witten associator and the latter is the non-trivial
associator. Indeed a normalized $2$-cochain $\beta$ satisfies
\[
 (\delta\beta)(\sigma,\sigma,\sigma)
 =
 \beta(\sigma,\sigma)\beta(1,\sigma)^{-1}
 \beta(\sigma,1)\beta(\sigma,\sigma)^{-1}
 =1,
\]
so the triple sign is not changed by normalized coboundary. Thus the
displayed sign is precisely the obstruction datum required by
Theorem~\ref{thm:monster-chain-level-e3-top}. The determinant
$2^{24}$ is only the determinant component of a local verification of
this datum.
\end{proof}

%%%%%%%%%%%%%%%%%%%%%%%%%%%%%%%%%%%%%%%%%%%%%%%%%%%%%%%%%%%%%%%%%%%%%%%%%%%%%
\section{Chain-level quasi-isomorphism}
\label{sec:monster-chain-qi}
%%%%%%%%%%%%%%%%%%%%%%%%%%%%%%%%%%%%%%%%%%%%%%%%%%%%%%%%%%%%%%%%%%%%%%%%%%%%%

Theorem~\ref{thm:monster-chain-level-e3-top} asserts a descended
chain-level $E_3$-topological structure only under the BV descent and
boundary-identification hypotheses. The boundary comparison is the
following conditional map.

\begin{proposition}[Chain-level FLM boundary comparison;
\ClaimStatusConditional]
\label{prop:monster-flm-chain-qi}
Let $\cF_\Lambda^{\mathrm{orb}} = \cF_\Lambda^{\mathrm{inv}} \oplus
\cF_\Lambda^{\mathrm{tw},+}$ denote the orbifold Costello-Gwilliam
BV factorisation algebra associated to abelian holomorphic Chern-
Simons on $X \times \R$ with charge lattice $\Lambda$ gauged by
$\Z/2$. Assume the finite-orbifold BV descent datum of
Theorem~\ref{thm:monster-chain-level-e3-top} and assume boundary
restriction identifies the invariant and twisted BV sectors with the
FLM sectors. Then there is a chain-level quasi-isomorphism
\[
 \Psi\colon V^\natural
 \;\xrightarrow{\sim}\;
 \cF_\Lambda^{\mathrm{orb}}
\]
of BV complexes, compatible with boundary restriction and intertwining
the FLM orbifold differential with the Costello-Gwilliam orbifold
BV differential.
\end{proposition}

\begin{proof}
The map $\Psi$ sends FLM invariant-sector generators
$\{v \otimes e^\lambda : \widetilde{\sigma} v = v,\ \sigma\lambda =
\lambda\}$ to their boundary-observable image in
$\cF_\Lambda^{\mathrm{inv}}$ under Costello-Li boundary restriction,
and FLM twisted-sector generators (elements of the unique irreducible
$\sigma$-twisted $V_\Lambda$-module) to their image in
$\cF_\Lambda^{\mathrm{tw},+}$. Three observations:
\begin{enumerate}[label=\textup{(\roman*)}]
\item $\Psi$ is a chain map because the FLM orbifold differential is
identified, by hypothesis, with the BRST differential restricted to
$\widetilde{\sigma}$-invariants and extended by the twisted sector.
\item Cohomology matches in each sector by the Costello-Li
abelian-CS boundary identification for lattice VOAs
(Theorem~\ref*{thm:E3-topological-km} restricted to abelian charge
lattice), together with the assumed boundary comparison for the
twisted sector.
\item Anomaly matching is exactly the local datum
$\alpha_{\mathrm{orb}}(\sigma,\sigma,\sigma)=+1$ from
Proposition~\ref{prop:monster-alpha-explicit-zero}; without this
datum the orbifold multiplication may carry a Dijkgraaf-Witten
associator.
\end{enumerate}
The graded character comparison with $J(\tau)$ is a VOA consistency
check, not a substitute for the BV anomaly trivialization.
\end{proof}

%%%%%%%%%%%%%%%%%%%%%%%%%%%%%%%%%%%%%%%%%%%%%%%%%%%%%%%%%%%%%%%%%%%%%%%%%%%%%
\section{Relation to Schellekens holomorphic $c=24$ VOAs}
\label{sec:monster-schellekens-extension}
%%%%%%%%%%%%%%%%%%%%%%%%%%%%%%%%%%%%%%%%%%%%%%%%%%%%%%%%%%%%%%%%%%%%%%%%%%%%%

\begin{remark}[Schellekens extension is a separate stratified theorem]
\label{rem:schellekens-71-honest}
The Monster criterion supplies the Type~B input in the Schellekens
$c=24$ landscape once the Leech $\Z/2$ BV anomaly datum is trivial.
The full $71$-entry statement is not a corollary of
$\Lambda^\sigma=0$; it is the separate three-stratum criterion
(Theorem~\ref{thm:schellekens-71-all-alpha-zero},
\texttt{schellekens\_71\_alpha\_classification\_platonic.tex}):

\begin{enumerate}[nosep]
\item \emph{Type A} (24 pure Niemeier lattice VOAs): trivial orbifold;
 abelian holomorphic Chern-Simons on the even-unimodular lattice.
\item \emph{Type B} (Leech $\mathbb{Z}/2 = V^{\natural}$, 1 case):
 $\Lambda^{\sigma} = 0$ because $\sigma=-1$ on a torsion-free
 lattice; Leech rootlessness identifies the moonshine entry, and
 the present chapter isolates the required local
 $\alpha_{\mathrm{DW}} = 0$ datum.
\item \emph{Type C} (Leech $\mathbb{Z}/n$, $n \in \{3,4,5,6,7,8,10,\geq 12\}$,
 46 cases): $\Lambda^{\sigma} \ne 0$ generally, but the
 van~Ekeren-Lam-M\"oller-Scheithauer 2021 (arXiv:2106.15491)
 level-matching condition supplies the cyclic-orbifold VOA input;
 the DW trivialization and finite-orbifold BV descent remain
 type-specific inputs.
\end{enumerate}

Total:
$24_{\mathrm{Niemeier}} + 1_{V_1=0} + 46_{\mathrm{nonlat}} = 71$,
with each stratum distinct in mechanism.
The bijection to generalised deep holes of Leech (M\"oller-Scheithauer
2023, arXiv:2303.11823) makes the Type~C enumeration canonical. This
chapter proves only the Type~B mechanism; Types~A and~C are imported
from the Schellekens classification chapter.

Citation chain for the cyclic-orbifold input:
\begin{itemize}[nosep]
\item \textbf{Cyclic-orbifold level matching} (van~Ekeren--M\"oller--
 Scheithauer, with the Schellekens-list application of
 van~Ekeren--Lam--M\"oller--Shimakura) supplies the VOA
 level-matching condition for the cyclic orbifold strata.
\item \textbf{M\"oller--Scheithauer} supplies the generalised-deep-hole
 enumeration of the nonzero $V_1$ Schellekens entries; this is the
 canonical source for the Type~C stratification, not the Monster
 argument.
\item \textbf{DLM 2000} (Dong-Li-Mason) provides the modular-invariance
 consistency for the resulting holomorphic orbifold VOA.
\end{itemize}

The Monster case (Type~B), the pure Niemeier cases (Type~A), and the
46 Type~C cases give chain-level $E_3$-topological descent only after
the stratum-specific anomaly and BV-descent data in
Theorem~\ref{thm:schellekens-71-all-alpha-zero} are supplied. The
Monster argument contributes one mechanism, not a uniform proof.
\end{remark}

%%%%%%%%%%%%%%%%%%%%%%%%%%%%%%%%%%%%%%%%%%%%%%%%%%%%%%%%%%%%%%%%%%%%%%%%%%%%%
\section{Summary}
\label{sec:monster-chain-level-summary}
%%%%%%%%%%%%%%%%%%%%%%%%%%%%%%%%%%%%%%%%%%%%%%%%%%%%%%%%%%%%%%%%%%%%%%%%%%%%%

\begin{itemize}
\item The DW anomaly class $\alpha_{\mathrm{orb}}$ of the Leech
$\Z/2$-orbifold is the finite-orbifold BV datum detected by
$\alpha_{\mathrm{orb}}(\sigma,\sigma,\sigma)=+1$
(Proposition~\ref{prop:monster-alpha-explicit-zero}).
\item The descended chain-level $E_3$-topological HT model with
boundary $V^\natural$ follows from
Theorem~\ref{thm:monster-chain-level-e3-top} only after the abelian
Leech HT model, anomaly trivialization, and finite-orbifold BV descent
are supplied.
\item The determinant $\det(1-\sigma|_\Lambda)=2^{24}$ contributes
the positive real determinant sign; it is not by itself the
Dijkgraaf-Witten computation.
\item The FLM twisted-sector Jacobi identity of
Frenkel-Lepowsky-Meurman~1988 supplies the chain-level associativity
of the boundary VOA product; the HT orbifold associator is controlled
by the BV anomaly datum.
\item Extension to the full Schellekens $c = 24$ list is handled by
the separate stratified theorem
\ref{thm:schellekens-71-all-alpha-zero}; the Monster argument supplies
only the Type~B summand because uniform $\Lambda^\sigma = 0$ does not
hold.
\end{itemize}
